% GENERATED FILE. Edit canonical lessons, then run npm run generate:books.
% canonical-source-hash: fnv1a64:50c98c4a7264c631
% canonical-lessons: ML-C423-nayam, ML-C423-natappakkuka, ML-C423-pasakkuka, ML-C423-bhedagati-ceyyuka, ML-C423-satyapratijna

\chapter{Making Policy}
\label{ch:ml-making-policy}
\hlchaptermodality{423}

\begin{chapteropening}
\textbf{By the end of this chapter:} \emph{I can say a policy, to implement, to pass a law, to amend and an oath.}

Complete the last lesson of chapter 423: I can say a policy, to implement, to pass a law, to amend and an oath.
\end{chapteropening}

\section[\texorpdfstring{\ml{നയം} (nayaṁ)}{നയം (nayaṁ)}]{\ml{നയം} (nayaṁ) — a policy}

\label{lesson:ML-C423-nayam}



\begin{quote}
Before the new one: say the Malayalam for to represent, then the Malayalam for the opposition.
\end{quote}

\subsection*{You'll want to know: \ml{നയം}}
\textbf{\ml{നയം}} — \emph{nayaṁ} — \textquotedblleft{}a policy\textquotedblright{}.

\begin{quote}
\ml{സർക്കാരിന്റെ} \ml{പുതിയ} \ml{നയം} \ml{കർഷകരെ} \ml{സഹായിക്കും}.
\end{quote}

\begin{quote}
\emph{sarkkārinṟe putiya nayaṁ karṣakare sahāyikkuṁ.}
\end{quote}

\begin{quote}
The government's new policy will help farmers.
\end{quote}

Plural: \ml{നയങ്ങൾ}.

\begin{cousinweb}
\textbf{\ml{നയം}} (\emph{nayaṁ}) is borrowed from Sanskrit \emph{naya}, \textquotedblleft{}conduct, a plan; wisdom\textquotedblright{} (CDIAL 6965), from the root \emph{nī}, \textquotedblleft{}to lead\textquotedblright{}.
\end{cousinweb}

\begin{grammarlens}[title={What you've built}]
One more word for giving reasons and opinions.
\end{grammarlens}

\subsection*{Your turn}
\begin{itemize}
\raggedright
  \item \emph{Say it:} \emph{nayaṁ}
  \item \emph{Say it:} \emph{nayaṁ}, once more
  \item \emph{Say it:} \emph{sarkkārinṟe putiya nayaṁ karṣakare sahāyikkuṁ.} — The government's new policy will help farmers.
  \item \emph{Recall:} say the Malayalam for to represent, then the Malayalam for the opposition, then say \emph{nayaṁ} again
\end{itemize}

\begin{tcolorbox}[breakable,colback=teal!4,colframe=teal!35!black,title={Before you move on}]
What does \ml{നയം} mean? (\textquotedblleft{}A policy\textquotedblright{}.) Say it once more.
\end{tcolorbox}
\section[\texorpdfstring{\ml{നടപ്പാക്കുക} (naṭappākkuka)}{നടപ്പാക്കുക (naṭappākkuka)}]{\ml{നടപ്പാക്കുക} (naṭappākkuka) — to implement}

\label{lesson:ML-C423-natappakkuka}



\begin{quote}
Before the new one: say the Malayalam for the opposition, then the Malayalam for a policy.
\end{quote}

\subsection*{You'll want to know: \ml{നടപ്പാക്കുക}}
\textbf{\ml{നടപ്പാക്കുക}} — \emph{naṭappākkuka} — \textquotedblleft{}to implement\textquotedblright{}.

\begin{quote}
\ml{പഞ്ചായത്ത്} \ml{പദ്ധതി} \ml{വേഗം} \ml{നടപ്പാക്കി}.
\end{quote}

\begin{quote}
\emph{pañcāyattŭ paddhati vēgaṁ naṭappākki.}
\end{quote}

\begin{quote}
The panchayat implemented the scheme quickly.
\end{quote}

Past: \ml{നടപ്പാക്കി}, \emph{naṭappākki}.

Also \ml{നടപ്പിലാക്കുക}, with the same meaning.

\begin{cousinweb}
\textbf{\ml{നടപ്പാക്കുക}} (\emph{naṭappākkuka}) is a native Malayalam compound: \ml{നടപ്പ്}, \textquotedblleft{}a going, a custom\textquotedblright{}, from \ml{നടക്കുക}, \textquotedblleft{}to walk; to be carried out\textquotedblright{} (DEDR 3582), and \ml{ആക്കുക}, \textquotedblleft{}to make\textquotedblright{} (DEDR 333): to put into practice.
\end{cousinweb}

\begin{grammarlens}[title={What you've built}]
One more word for giving reasons and opinions.
\end{grammarlens}

\subsection*{Your turn}
\begin{itemize}
\raggedright
  \item \emph{Say it:} \emph{naṭappākkuka}
  \item \emph{Say it:} \emph{naṭappākkuka}, once more
  \item \emph{Say it:} \emph{pañcāyattŭ paddhati vēgaṁ naṭappākki.} — The panchayat implemented the scheme quickly.
  \item \emph{Recall:} say the Malayalam for the opposition, then the Malayalam for a policy, then say \emph{naṭappākkuka} again
\end{itemize}

\begin{tcolorbox}[breakable,colback=teal!4,colframe=teal!35!black,title={Before you move on}]
What does \ml{നടപ്പാക്കുക} mean? (\textquotedblleft{}To implement\textquotedblright{}.) Say it once more.
\end{tcolorbox}
\section[\texorpdfstring{\ml{പാസാക്കുക} (pāsākkuka)}{പാസാക്കുക (pāsākkuka)}]{\ml{പാസാക്കുക} (pāsākkuka) — to pass a law; to approve}

\label{lesson:ML-C423-pasakkuka}



\begin{quote}
Before the new one: say the Malayalam for a policy, then the Malayalam for to implement.
\end{quote}

\subsection*{You'll want to know: \ml{പാസാക്കുക}}
\textbf{\ml{പാസാക്കുക}} — \emph{pāsākkuka} — \textquotedblleft{}to pass a law; to approve\textquotedblright{}.

\begin{quote}
\ml{നിയമസഭ} \ml{പുതിയ} \ml{നിയമം} \ml{പാസാക്കി}.
\end{quote}

\begin{quote}
\emph{niyamasabha putiya niyamaṁ pāsākki.}
\end{quote}

\begin{quote}
The assembly passed a new law.
\end{quote}

Past: \ml{പാസാക്കി}, \emph{pāsākki}.

\ml{ആക്കുക}, to make, turns a loan into a verb.

\begin{cousinweb}
\textbf{\ml{പാസാക്കുക}} (\emph{pāsākkuka}) is a Malayalam compound: the first part was borrowed from English \emph{pass}, and \ml{ആക്കുക}, \textquotedblleft{}to make\textquotedblright{}, is native Dravidian (DEDR 333).
\end{cousinweb}

\begin{grammarlens}[title={What you've built}]
One more word for giving reasons and opinions.
\end{grammarlens}

\subsection*{Your turn}
\begin{itemize}
\raggedright
  \item \emph{Say it:} \emph{pāsākkuka}
  \item \emph{Say it:} \emph{pāsākkuka}, once more
  \item \emph{Say it:} \emph{niyamasabha putiya niyamaṁ pāsākki.} — The assembly passed a new law.
  \item \emph{Recall:} say the Malayalam for a policy, then the Malayalam for to implement, then say \emph{pāsākkuka} again
\end{itemize}

\begin{tcolorbox}[breakable,colback=teal!4,colframe=teal!35!black,title={Before you move on}]
What does \ml{പാസാക്കുക} mean? (\textquotedblleft{}To pass a law; to approve\textquotedblright{}.) Say it once more.
\end{tcolorbox}
\section[\texorpdfstring{\ml{ഭേദഗതി} \ml{ചെയ്യുക} (bhēdagati ceyyuka)}{ഭേദഗതി ചെയ്യുക (bhēdagati ceyyuka)}]{\ml{ഭേദഗതി} \ml{ചെയ്യുക} (bhēdagati ceyyuka) — to amend}

\label{lesson:ML-C423-bhedagati-ceyyuka}



\begin{quote}
Before the new one: say the Malayalam for to implement, then the Malayalam for to pass a law.
\end{quote}

\subsection*{You'll want to know: \ml{ഭേദഗതി} \ml{ചെയ്യുക}}
\textbf{\ml{ഭേദഗതി} \ml{ചെയ്യുക}} — \emph{bhēdagati ceyyuka} — \textquotedblleft{}to amend\textquotedblright{}.

\begin{quote}
\ml{സർക്കാർ} \ml{ആ} \ml{നിയമം} \ml{ഭേദഗതി} \ml{ചെയ്തു}.
\end{quote}

\begin{quote}
\emph{sarkkār ā niyamaṁ bhēdagati ceytu.}
\end{quote}

\begin{quote}
The government amended that law.
\end{quote}

Past: \ml{ഭേദഗതി} \ml{ചെയ്തു}, \emph{bhēdagati ceytu}.

The noun \ml{ഭേദഗതി} is an amendment.

\begin{cousinweb}
\textbf{\ml{ഭേദഗതി} \ml{ചെയ്യുക}} (\emph{bhēdagati ceyyuka}) is a phrase verb built on the noun \ml{ഭേദഗതി}, \textquotedblleft{}an amendment\textquotedblright{}, a Malayalam compound of Sanskrit \emph{bheda}, \textquotedblleft{}breaking, a change\textquotedblright{} (Monier-Williams), and \emph{gati}, \textquotedblleft{}going\textquotedblright{} (CDIAL 4009); \ml{ചെയ്യുക}, \textquotedblleft{}to do\textquotedblright{}, makes it a verb.
\end{cousinweb}

\begin{grammarlens}[title={What you've built}]
One more word for giving reasons and opinions.
\end{grammarlens}

\subsection*{Your turn}
\begin{itemize}
\raggedright
  \item \emph{Say it:} \emph{bhēdagati ceyyuka}
  \item \emph{Say it:} \emph{bhēdagati ceyyuka}, once more
  \item \emph{Say it:} \emph{sarkkār ā niyamaṁ bhēdagati ceytu.} — The government amended that law.
  \item \emph{Recall:} say the Malayalam for to implement, then the Malayalam for to pass a law, then say \emph{bhēdagati ceyyuka} again
\end{itemize}

\begin{tcolorbox}[breakable,colback=teal!4,colframe=teal!35!black,title={Before you move on}]
What does \ml{ഭേദഗതി} \ml{ചെയ്യുക} mean? (\textquotedblleft{}To amend\textquotedblright{}.) Say it once more.
\end{tcolorbox}
\section[\texorpdfstring{\ml{സത്യപ്രതിജ്ഞ} (satyapratijña)}{സത്യപ്രതിജ്ഞ (satyapratijña)}]{\ml{സത്യപ്രതിജ്ഞ} (satyapratijña) — an oath}

\label{lesson:ML-C423-satyapratijna}



\begin{quote}
Before the new one: say the Malayalam for to pass a law, then the Malayalam for to amend.
\end{quote}

\subsection*{You'll want to know: \ml{സത്യപ്രതിജ്ഞ}}
\textbf{\ml{സത്യപ്രതിജ്ഞ}} — \emph{satyapratijña} — \textquotedblleft{}an oath\textquotedblright{}.

\begin{quote}
\ml{പുതിയ} \ml{മന്ത്രിമാർ} \ml{ഇന്ന്} \ml{സത്യപ്രതിജ്ഞ} \ml{ചെയ്തു}.
\end{quote}

\begin{quote}
\emph{putiya mantrimār innŭ satyapratijña ceytu.}
\end{quote}

\begin{quote}
The new ministers took the oath today.
\end{quote}

To take an oath is \ml{സത്യപ്രതിജ്ഞ} \ml{ചെയ്യുക}.

\begin{cousinweb}
\textbf{\ml{സത്യപ്രതിജ്ഞ}} (\emph{satyapratijña}) is built from Sanskrit \emph{satya}, \textquotedblleft{}true; truth\textquotedblright{} (CDIAL 13112), and \emph{pratijñā}, \textquotedblleft{}a promise\textquotedblright{} (CDIAL 8563): a promise of truth.
\end{cousinweb}

\begin{grammarlens}[title={What you've built}]
One more word for giving reasons and opinions.
\end{grammarlens}

\subsection*{Your turn}
\begin{itemize}
\raggedright
  \item \emph{Say it:} \emph{satyapratijña}
  \item \emph{Say it:} \emph{satyapratijña}, once more
  \item \emph{Say it:} \emph{putiya mantrimār innŭ satyapratijña ceytu.} — The new ministers took the oath today.
  \item \emph{Recall:} say the Malayalam for to pass a law, then the Malayalam for to amend, then say \emph{satyapratijña} again
\end{itemize}

\begin{tcolorbox}[breakable,colback=teal!4,colframe=teal!35!black,title={Before you move on}]
What does \ml{സത്യപ്രതിജ്ഞ} mean? (\textquotedblleft{}An oath\textquotedblright{}.) Say it once more.
\end{tcolorbox}
